\documentclass[11pt]{article}

\usepackage[T1]{fontenc}
\usepackage[utf8]{inputenc}
\usepackage[margin=1in]{geometry}
\usepackage{amsmath,amssymb}
\usepackage{graphicx}
\usepackage{booktabs}
\usepackage{longtable}
\usepackage{enumitem}
\usepackage{caption}
\usepackage[colorlinks=true,linkcolor=blue,citecolor=blue,urlcolor=blue]{hyperref}
\usepackage{titlesec}

\captionsetup{font=small,labelfont=bf}
\setlist{itemsep=2pt,parsep=0pt,topsep=3pt}
\titlespacing*{\section}{0pt}{2.2ex}{1.0ex}
\titlespacing*{\subsection}{0pt}{1.8ex}{0.7ex}

% Looks in figures/ first, then the project root, so either layout works.
\newcommand{\rfig}[4][0.95]{%
  \begin{figure}[htbp]\centering
  \IfFileExists{figures/#2}%
    {\includegraphics[width=#1\textwidth]{figures/#2}}%
    {\IfFileExists{#2}%
      {\includegraphics[width=#1\textwidth]{#2}}%
      {\fbox{\parbox{0.85\textwidth}{\centering\vspace{2cm}\small\itshape
       missing: \texttt{\detokenize{#2}}\vspace{2cm}}}}}%
  \caption{#4}\label{fig:#3}\end{figure}}

\title{\vspace{-2em}A Continuous Attractor in Mouse Motor Cortex,\\
and a Causal Readout for Closed-Loop Intervention}
\author{
Samridh Tiwari$^{1,2}$ \quad Om Kolhe$^{1}$ \quad Krishna Jayant$^{1}$\\[4pt]
\small $^{1}$Nano Neurotechnology Laboratory, Weldon School of Biomedical
Engineering,\\
\small Purdue University, West Lafayette, IN 47907\\
\small $^{2}$Indian Institute of Technology Bombay, Mumbai 400076, India\\[4pt]
\small \texttt{samridhtiwari3005@gmail.com}}
\date{SURF 2026 Final Report}

\begin{document}
\maketitle

\begin{abstract}\noindent
A motor plan has to do two things that pull against each other. It has to stay
put, so that noise does not corrupt it, and it has to be movable, so that a cue
can change it. What kind of circuit does both has not been shown in mouse
forelimb motor cortex.

We fit a dynamical model, LFADS, to spiking from two mice (9 sessions, 39 to 118
neurons recorded at once) during a cued lever-pull task, then treated the fitted
model as a dynamical system and examined its shape.

The learned state passes three tests. It predicts neurons the model never saw
better than a standard baseline, in 8 of 8 sessions, by a median factor of
$4.1\times$. A linear readout of it recovers the \emph{actual lever movement} on
held-out trials at $R^2 = 0.47$ to $0.77$, even though the model was fit to
spikes and never shown the lever. And starting from one inferred state with all
external input switched off, the learned rule regenerates 94\% of a single
trial's latent trajectory, against 2\% for a predictor that just uses the
average trial.

Examining the rule at 512 of its fixed points reveals a \textbf{continuous
attractor}: about 33 directions in which a push stays put, about 95 in which a
push dies away within 84~ms, and essentially none in which it grows. The same
structure appears in five separately trained models per animal and in both
animals. Two models trained independently on different mice produce the same
trajectory shape once rigidly aligned (shape difference 0.11, where 1.00 is
chance). Given only spikes, the model also infers an external input that peaks
100 to 140~ms after the cue and is larger on trials the mouse got right. A push
along the manifold keeps 75\% of its size after 1.2~s; a push across it keeps
8\%.

Every one of those estimates uses information from the future, so none can run
live. We therefore built a \textbf{causal real-time detector} of impending
movement from unsorted multiunit activity. Streaming from the amplifier it
flagged movement a median of \textbf{80~ms before the lever moved}, with 90\% of
alarms early and 71\% followed by a real movement, using \textbf{under 1~ms of
compute per 20~ms bin}. Because the animal moves every 3.7~s, we scored it
against a null that preserves alarm rate; precision beat chance by 6.7 standard
deviations. Applied to the next day's recording it dropped to chance (AUC 0.485)
and recovered to 0.740 after a rigid latent alignment, against a same-day
ceiling of 0.768, which shows the readout coordinates rotate while the dynamics
do not.

We also report two lessons about method. Choosing a model by reconstruction
error picks models that explain the data with injected input rather than with
their own dynamics, and those models are useless for studying dynamics. And for
live use, a bounded score should be preferred to an unbounded one even when
offline accuracy is identical.
\end{abstract}

\tableofcontents
\bigskip

\section{Introduction and Background}

Before you move, motor cortex is already busy. Activity starts changing hundreds
of milliseconds before the limb does, and during that period it is not driving
muscles \cite{tanji1976,churchland2012}. Population recordings show this activity
lives in a small number of dimensions and traces repeatable paths
\cite{cunningham2014}, and that the path, rather than any single neuron's tuning,
is the right thing to describe \cite{churchland2012,vyas2020}.

That sets up a problem. The plan has to be \emph{stable}, or noise would corrupt
it before the movement starts. But it also has to be \emph{changeable}, or a cue
could never move it. A network that resists everything cannot be updated; one
that yields to everything cannot hold anything.

A classical answer is a \textbf{continuous attractor}: a connected surface of
states, each of which the network is content to sit in
\cite{seung1996,zhang1996}. Push the state along that surface and it stays where
you put it, so the surface stores information. Push it off the surface and the
network drags it back, so noise is rejected. Continuous attractors are
established for head direction and for holding the eyes still
\cite{zhang1996,chaudhuri2019}, but evidence in motor preparation is thinner,
mainly because finding one requires a model that \emph{generates} the dynamics
rather than a description of the paths taken.

LFADS \cite{pandarinath2018} is such a model. It is an autoencoder whose decoder
is a recurrent network trained to produce the observed spike trains. Once fit,
that recurrent network is a dynamical system you can interrogate directly: find
where it holds still, compute exactly how it responds to a small push, and read
stability off the result \cite{sussillo2013,golub2018}.

Our reason for wanting this is practical. To nudge a motor plan back toward a
working configuration you need three things: somewhere to aim, a model of what a
push actually does, and a live estimate of where the circuit is right now. Part~A
of this report supplies the first two. Part~B works on the third and says exactly
what is still missing.

\paragraph{What this report contributes.}
\begin{enumerate}
\item Evidence for a continuous attractor in mouse forelimb motor cortex,
reproduced in five separately trained models per animal and in two animals.
\item A model fit only to spikes that nonetheless predicts the continuous lever
movement on trials it never saw.
\item A cue-locked external input recovered from spikes alone, without the model
ever being told when the cue happened or which trials succeeded.
\item A measurement of how the cost of intervening depends on which direction you
push, and where the behaviourally useful direction sits relative to that.
\item A causal movement detector, validated live against a null that corrects for
how often the animal moves anyway, with its cross-day failure characterised and
fixed.
\item Two methodological warnings, one about choosing models and one about
choosing readouts for live use.
\end{enumerate}

\section{Methods}

\subsection{Task and recordings}

Two mice (\texttt{61199SOM}, \texttt{61198SOM}) did a cued forelimb lever-pull
task. A tone marked the start of a trial. If the mouse pulled the lever properly
it got water (``Hit''); if it did not pull, it got nothing (``Miss''). Both
outcomes happen often enough to compare.

We recorded spiking from forelimb motor cortex across \textbf{9 sessions}, four
in one animal and five in the other, with \textbf{39 to 118} neurons isolated
simultaneously per session. Trials were cut from 1.5~s before the tone to 1.5~s
after it. LFADS was fit on the 8 sessions present in the stitched archive:
\texttt{61199SOM}, 1163 trials (844 Hit); \texttt{61198SOM}, 956 trials (595
Hit).

Part~B used further sessions recorded with a 64-channel silicon probe at 20~kHz
(Intan RHD2000), with the lever voltage and the cue and reward TTL lines on the
same clock so everything is directly comparable. The main reference session has
280 cues, 206 pulls and 197 rewards over about 37.5~minutes. Two further sessions
recorded primary (M1) and premotor (M2) cortex at the same time.

\subsection{Preprocessing}

For Part~A, 1~ms spike counts were summed into \textbf{20~ms bins}, giving 150
bins per trial. Trials were split 80/20 into training and validation by time, not
at random. Nothing was smoothed before fitting; smoothing ($\sigma = 50$~ms) was
used only for the descriptive plots in \S\ref{sec:lowd}.

For Part~B \textbf{no spike sorting was done at any stage}. The detector counts
threshold crossings instead. This keeps the population structure and removes a
step that cannot be done reliably while a session is running
\cite{trautmann2019}.

\subsection{The model}

\textbf{LFADS} \cite{pandarinath2018} was fit with \texttt{lfads-torch}
\cite{sedler2023}. Several sessions from one animal were fit together, with a
separate input and output matrix per session but one shared generator, so the
latent variables mean the same thing across that animal's sessions
\cite{keshtkaran2022}.

The pieces, in the order data flows through them. Spikes go in. An encoder
compresses each trial into an \emph{initial state}. The \emph{generator}, a
128-unit recurrent network, carries that state forward in 20~ms steps. Its state
at each step is turned into 15 \emph{factors}, and those factors give a firing
rate for every neuron. Those predicted rates are compared against the real
spikes, and that comparison is the only thing training optimises. A separate
\emph{controller} supplies \emph{inferred inputs} back into the generator, which
is the model's account of what arrived from outside on that trial.

Nothing in this encodes dynamics, or the cue, or which trials succeeded.
Everything reported below is therefore something the model found, not something
we put in.

Other settings: controller output dimension 4, initial-condition dimension 64,
dropout 0.1, batch size 32, learning rate $4\times10^{-3}$, Poisson spike model.
\textbf{GPFA} \cite{yu2009} was fit at the \emph{same} latent dimensionality
(\texttt{xDim} = 15) on the same sessions, as the comparison.

\paragraph{A note on which quantity is used where.} The generator's internal
state is 128-dimensional; the factors are 15-dimensional. The lever decoding
(\S\ref{sec:lever}), the flow field (\S\ref{sec:flow}) and the fixed-point
analysis (\S\ref{sec:attractor}) all use the \emph{generator state}. Only the
cross-animal comparison (\S\ref{sec:gen}) uses the \emph{factors}.

\subsection{Choosing which fitted model to use}\label{sec:selmethod}

A sweep over regularisation settings produced 10 fitted models per animal. The
obvious way to pick one is validation reconstruction error. We also measured
something else, which we call the \textbf{autonomy score}: run the generator
forward from its inferred initial state with all inferred input set to zero, and
correlate that trajectory with the one the full model produced. A high score
means the learned rule is carrying the trial. A low score means the inferred
inputs are doing the work and the rule is mostly decorative.

We selected \texttt{cfg06} for both animals, which is \emph{not} the
reconstruction winner, for reasons shown in \S\ref{sec:sel}. This is a deliberate
departure from standard practice (\S\ref{sec:dev}, item 1).

\subsection{How the model was tested}

\paragraph{Predicting neurons it never saw \cite{pei2021}.} Split the neurons
into 5 groups and the trials into 5 groups. Hide one group of neurons entirely,
infer the latent state from the neurons that remain, then use that state to
predict the hidden neurons on held-out trials. Score in bits per spike relative
to each neuron's own average rate. GPFA and LFADS were matched to identical
trials, scored against an identical spike target, and read out through an
identical Poisson GLM, so what differs is the quality of the latent state and
nothing else.

\paragraph{Predicting the lever.} Lever position at time $t$ was predicted from a
\emph{window} of generator state spanning $-60$ to $+160$~ms around $t$, by ridge
regression. Motor cortex leads movement, so predicting from the single current
bin is simply the wrong model: it scores $R^2 = 0.32$ on data where the windowed
version reaches $0.77$. Cross-validation was 5-fold across trials, with an inner
3-fold loop choosing the ridge penalty. The comparison is against the
training-set average pull applied to every held-out trial. Miss trials were left
out of scoring because the lever barely moves on them, which makes $R^2$
meaningless.

\paragraph{Letting the rule run on its own.} Start from a trial's inferred
initial state, set all input to zero, and iterate the generator 150 steps. Score
how much of that trial's real latent trajectory is recovered. Three controls:
shuffle which trial's starting state is used, compare against a
condition-average predictor, and check training against held-out trials.

\subsection{Examining the fitted dynamics}

\paragraph{Finding where the network holds still.} We searched for states where
the rule leaves the state unchanged, by minimising
$\|f(\mathbf{g},\mathbf{0}) - \mathbf{g}\|^2$ with Adam (learning rate
$2\times10^{-2}$, cosine annealing to $10^{-5}$, 12{,}000 iterations) from 512
starting points drawn from real recorded states with noise added ($\sigma =
0.05$) \cite{sussillo2013,golub2018}. A candidate was kept if its relative
displacement fell below $10^{-3}$.

\paragraph{Reading off stability.} At each fixed point we computed the Jacobian
exactly by automatic differentiation and took its eigenvalues. A direction is
\textbf{marginal} if $\big||\lambda|-1\big| < 0.01$ (a push neither grows nor
shrinks), \textbf{stable} if $|\lambda| < 0.99$ (a push decays), and
\textbf{unstable} if $|\lambda| > 1.01$ (a push grows). Decay time is
$\tau = -\Delta t / \ln|\lambda|$.

\paragraph{The flow field.} At each point on a grid we took the \emph{average
real 128-dimensional state} of nearby data, advanced it one step through the
trained network, and projected the change back to two dimensions. Every arrow is
therefore a real one-step prediction at a place the recordings actually visit,
not an extrapolation into empty space. Arrows fade where less data supports them.

\paragraph{Pushing the network.} We displaced the full nonlinear generator from a
fixed point by a unit-length vector, either along a marginal direction or along a
randomly chosen non-marginal one, released it, and watched how much displacement
survived over 1.2~s. The eigen-decomposition only chooses the direction; the
simulation uses the real trained network, not its linear approximation.

\subsection{The real-time detector}

Everything above depends on the inferred initial state, which LFADS computes
using an encoder that reads the \emph{whole} trial, including what happens after
the moment being estimated (\S\ref{sec:tips}, tip 4). None of it can run live.
Part~B therefore asks a smaller question: can we get the \emph{trigger} a
closed-loop system needs, causally, and fast enough?

\paragraph{The pipeline.} One pass over the raw stream. Subtract the median
across the 64 channels. Apply a causal 3rd-order Butterworth high-pass at 300~Hz.
Set a per-channel threshold at $4\times$ a running median-absolute-deviation
estimate (5~s window, 1~s hop). Count threshold crossings with a 1~ms refractory
period, accumulate them into \textbf{20~ms bins}, z-score each channel against
its running baseline, and stack the \textbf{last 12 bins (240~ms)} into a
$64 \times 12 = 768$-number feature vector. Filter and baseline state carry
across file boundaries and are never reset, which is why replaying a recording
gives exactly what the live system gave.

\paragraph{The decision.} Project that vector onto one direction and raise an
alarm when the score stays above threshold for two consecutive bins (40~ms). We
compared two projections: a regularised Fisher discriminant \cite{fisher1936},
and a PCA-whitened cosine similarity to a movement template, which is bounded to
$[-1,1]$ by construction.

\paragraph{Labels.} Training uses only the lever. Cue timing is never supplied,
so the detector cannot know whether a movement was cued or self-initiated.
Positives are the bins in the 200~ms before a pull that crossed the hardware
reward threshold (\S\ref{sec:tips}, tip 7). Negatives are bins more than 2~s from
any pull with the lever at rest. Cross-validation is \textbf{blocked in time},
because random splits leak between neighbouring trials.

\paragraph{What it is compared against.} An unsupervised CUSUM \cite{page1954}
ramp detector on the first principal component of the same features, which is the
standard approach and the floor. And a \textbf{circular-shift null}: rotate all
the alarm times around the session clock, which keeps the alarm rate and any
rhythm but destroys the alignment to movement. This is essential, because the
animal pulls every 3.7~s or so and a detector firing at random would still land
near movements often. Compute time was measured per bin on one CPU core; no GPU
is used at run time.

\subsection{Across days and across areas}

Cross-day transfer was tested by applying a model trained on one day directly to
the next, then repeating after aligning the two latent spaces by Procrustes
registration using pull-triggered averages as landmarks
\cite{gallego2020,degenhart2020}. For the M1 versus M2 comparison both probes
went through the identical front end and the same 20~ms binning, with matched
channel counts, time-blocked cross-validation, label-shuffle nulls, and bootstrap
intervals on the difference.

\subsection{Sample, Code, and Data Storage}

Analysis code, trained checkpoints, the derived numbers behind every reported
value, and the editable source of every figure are at\\
\url{https://github.com/Samridh-Tiwari/surf2026-motor-dynamics}.

Every figure is produced by a committed script that also writes a
\texttt{*\_stats.json} of the values it reports, and the \texttt{README} maps
each reported number to the script that produces it. Paths come from environment
variables, so nothing inside the scripts needs editing.

\textbf{Naming.} Animals are \texttt{61199SOM} and \texttt{61198SOM}. Sessions
are \texttt{Day2} to \texttt{Day7} within an animal. The stitched archive indexes
the same sessions as \texttt{dataset\_01} to \texttt{dataset\_08}. Fitted models
are \texttt{cfg01} to \texttt{cfg10}, and \texttt{cfg06} is the one used
throughout.

\textbf{Data.} The recordings are too large to host in the repository: the
stitched spike archive is 1.7~GB, the per-session behavioural files total 33~GB,
and the LFADS per-session outputs 528~MB. They are \textbf{provided on request}.
The committed checkpoints, 2.3~MB each, are enough to regenerate every Part~A
result without retraining.

\begingroup\small
\begin{longtable}{@{}p{0.44\textwidth}p{0.52\textwidth}@{}}
\caption{Where each reported number comes from.}\label{tab:prov}\\
\toprule
\textbf{Value} & \textbf{Script and output} \\
\midrule
\endfirsthead
\toprule
\textbf{Value} & \textbf{Script and output} \\
\midrule
\endhead
8 of 8 sessions, 0.303 vs 0.080 bits/spike, median $4.1\times$ &
\texttt{dyn\_recompute\_cosmooth.py} $\to$ \texttt{cosmooth\_sweep.json} \\
Lever decode $R^2$ 0.47--0.77, baseline 0.26--0.68 &
\texttt{fig\_lever\_decode.py} $\to$ \texttt{lever\_decode\_stats.json} \\
Reconstruction and autonomy for every fitted model &
\texttt{fig\_model\_selection.py} $\to$ \texttt{model\_selection\_stats.json} \\
Rule run alone: 94\% pre-cue, 82\% post-cue &
\texttt{dyn\_slide\_a7\_bridge.py} $\to$ \texttt{rollout\_stats.json} \\
512 fixed points, 33 marginal, 95 stable, $\tau$ 84~ms &
\texttt{fig3\_attractor\_cfg06.py} $\to$ \texttt{fig3\_attractor\_stats.json} \\
Marginal counts, 5 models $\times$ 2 animals &
\texttt{dyn\_robust\_61198.py} $\to$ \texttt{robustness*.json} \\
Flow field; 2 components hold 36\% and 59\% &
\texttt{fig\_flowfield.py} $\to$ \texttt{flowfield\_stats.json} \\
Across sessions 0.02--0.21, across animals 0.11 &
\texttt{fig\_generalize2.py} $\to$ \texttt{generalization\_stats.json} \\
Cue-locked input, $1.73\times$ and $2.18\times$ baseline &
\texttt{dyn\_slide\_a9.py} $\to$ \texttt{cue\_stats\_cfg06.json} \\
Push retention 75\% vs 8\%; $1.38\times$ chance &
\texttt{dyn\_slide\_a10.py} $\to$ \texttt{steering\_stats.json} \\
Dimensionality; Spearman $r = -0.16$ &
\texttt{reproduce\_lfads\_paper.py} $\to$ \texttt{dimensionality\_*.json} \\
Detector AUC 0.877; 60\% at $\sim$5 false alarms/min &
\texttt{m1\_realtime\_test.m} \\
Live: 80~ms lead, 90\% early, 71\% precision &
\texttt{closedloop\_live\_gui.m} $\to$ \texttt{run\_stats.mat} \\
Circular-shift null, $z = +6.7$ and $+8.2$ &
\texttt{shuffle\_control.m} \\
Cross-day 0.485 $\to$ 0.740, ceiling 0.768 &
\texttt{m1\_crossday.m}, \texttt{m1\_deploy\_crossday.m} \\
M1 vs M2, 0~ms lag, $r = 0.98$--$0.99$ &
\texttt{m1m2\_timing.m}, \texttt{m1m2\_cuefree\_test.m} \\
Reward threshold 2.8197 / 2.8188~V &
\texttt{find\_reward\_threshold.m}, \texttt{genuine\_pulls.m} \\
\bottomrule
\end{longtable}
\endgroup

Environment: conda \texttt{lfads-torch} with PyTorch and CUDA for Part~A, MATLAB
R2026a for Part~B. Everything random is seeded, so the reported numbers reproduce
exactly.

\subsection{Reference Methods}

LFADS: \url{https://github.com/arsedler9/lfads-torch} \cite{sedler2023}. GPFA:
\url{https://github.com/NeuralEnsemble/elephant} and the original release
\cite{yu2009}. Finding fixed points follows Sussillo and Barak
\cite{sussillo2013} and Golub and Sussillo \cite{golub2018}. The held-out-neuron
test follows the Neural Latents Benchmark protocol \cite{pei2021}. Working
without spike sorting follows Trautmann et al.\ \cite{trautmann2019}. Aligning
latent spaces across days follows Gallego et al.\ \cite{gallego2020} and
Degenhart et al.\ \cite{degenhart2020}.

\subsection{Insider Tips}\label{sec:tips}

Things that cost time and are not in any documentation.

\begin{enumerate}
\item The LFADS factor read-out is \texttt{KernelNormalizedLinear}, which
normalises each weight row before using it. Multiplying by the stored weights
does \textbf{not} give you the saved factors. Apply
\texttt{F.normalize(W, p=2, dim=1)} first.
\item \texttt{gen\_init} in the output file is the 128-dimensional generator
state \emph{after} the \texttt{ic\_to\_g0} projection, not the 64-dimensional
sample.
\item Saved \texttt{gen\_states} are averages over 50 posterior samples. Because
the network is nonlinear, re-simulating from the average starting state does not
reproduce them exactly (correlation about 0.99 over 150 steps). That is expected,
not a bug.
\item The encoder that produces the initial state reads the \textbf{whole trial,
both directions}. So anything computed from it, or from the factors, uses
information from the future. This matters for interpretation and is the whole
reason Part~B exists.
\item GPFA \texttt{.mat} files are MATLAB v7.3, which is HDF5 with struct arrays.
Fields have to be read by chasing object references.
\item Finding fixed points is sensitive to the learning-rate schedule. Cosine
annealing and a \emph{relative} convergence tolerance were both needed. The count
of marginal directions also shifts slightly with how many fixed points you sample
(33.2 at 60 sampled, 36.7 at 40), so always report the sample size next to the
count.
\item \textbf{The lever voltage that triggers reward is fixed in hardware, so
recover it as an absolute voltage, not as a statistic of the session.} We
measured 2.8197~V and 2.8188~V on two days, 0.9~mV apart. Thresholds based on
median absolute deviation or on amplitude percentiles drift between sessions and
quietly let tiny movements through: on one session that inflated the pull count
from 864 to 2702. Recover it from the lower edge of the peak lever voltage just
before each reward.
\item Behavioural events are in the raw \texttt{.rhd} files
(\texttt{ANALOG-IN-2} is the lever, \texttt{DIGITAL-IN-04} the cue,
\texttt{DIGITAL-IN-03} the reward), not in the exported workspace. The default
Intan reader throws away the analog board channels and has to be extended.
\item Carry the causal filter and baseline state across file boundaries.
Re-initialising per file creates a transient at every boundary that the detector
happily reports as movement.
\item Decoding movement from neural activity needs a \textbf{lagged} readout.
Using only the current bin scores $R^2 = 0.32$ where a window reaches $0.77$ on
the same data, purely because motor cortex leads movement.
\item Compare trajectory shapes in the model's \emph{full} latent space, not in a
2-D picture of it. The same comparison looks like a null result in two dimensions
holding 36\% of the variance, and is clean in 15 dimensions.
\end{enumerate}

\subsection{Deviations from Standard Methodology}\label{sec:dev}

\begin{enumerate}
\item \textbf{Model selection} on autonomy, subject to near-best reconstruction,
because reconstruction cannot see the difference that matters here
(\S\ref{sec:sel}).
\item \textbf{How we perturb.} We first tried time-varying LQR control through
the model's 4 inferred input channels. On the selected model that gives 0.1\%
error reduction, because those 4 channels have almost no authority over a
128-dimensional state. We therefore perturb the state directly, which is closer
to what a physical intervention would do.
\item \textbf{GPFA dimensionality} was fixed at 15 to match LFADS rather than
cross-validated, since the stored \texttt{cvLL} field is empty. This is a matched
comparison, not a comparison of each method at its own best setting.
\item \textbf{Readout choice} favours the bounded cosine over the Fisher
discriminant despite identical held-out accuracy, for robustness
(\S\ref{sec:readout}).
\item \textbf{Cross-validation} is blocked in time, not random. Random folds leak
between neighbouring trials and inflated one cross-area difference by about 40\%.
\end{enumerate}

\subsection{Methods Planned for Completion}

Two things are in progress. \textbf{Causal state estimation}: replace the
bidirectional encoder with a forward-only one over a fixed window of the past, so
the output is a position on the manifold rather than a yes/no alarm. And an
\textbf{optogenetic protocol} that stimulates along versus across the manifold,
triggered by the validated detector, to test \S\ref{sec:steer} directly.

\section{Results}

\subsection{The population uses few dimensions, and not because we recorded few
neurons}\label{sec:lowd}

About 3 principal components capture 95\% of the variance in trial-averaged
activity. The obvious worry is that this is an artefact of recording too few
cells. It is not: the number of components needed does not grow with the number
of neurons recorded, across a range from 39 to 118 (Spearman $r = -0.16$,
$p = 0.70$ at the 99\% criterion; Fig.~\ref{fig:dim}).

Smoothing does inflate apparent low-dimensionality, so we treat this as
justification for choosing a 15-dimensional latent space with room to spare,
rather than as a result in itself. The dimensionality result we actually rely on
comes from the dynamics (\S\ref{sec:attractor}).

\rfig[0.95]{dimensionality.png}{dim}{How many dimensions the population needs.
\emph{Left}, cumulative variance explained, one line per session, shaded by how
many neurons that session recorded. About 3 components reach 95\%. The dashed
line marks the 15 dimensions given to the model, chosen with headroom rather than
fitted. \emph{Right}, components needed against neurons recorded. If low
dimensionality were an undersampling artefact, these would trend upward. They do
not.}

\subsection{The latent state predicts neurons the model never saw}
\label{sec:cosmooth}

A model with enough parameters can fit anything, so the test has to be one that
memorisation cannot pass. We hid entire neurons, inferred the state from the
neurons that remained, and asked the model to predict the hidden ones.

LFADS won in \textbf{8 of 8 fitted sessions} across both animals (0.303 versus
0.080 bits per spike; median ratio $\mathbf{4.1\times}$, range 2.6 to
5.9$\times$; Fig.~\ref{fig:cosmooth}).

A Wilcoxon signed-rank test gives $p = 0.008$, and that number should be read
carefully. With 8 paired samples, $0.008$ is the \emph{smallest value the test
can produce} ($2/2^{8} = 0.0078$). It tells you every session went the same way
and nothing more. The size of the effect is the $4.1\times$; the consistency is
the 8 of 8.

We deliberately hold out \emph{neurons} rather than \emph{trials}. LFADS infers a
separate starting state and separate inputs for every trial, so it fits
individual trials almost by construction, which makes held-out-trial performance
a weak test.

\rfig[0.88]{fig1_model_validation.png}{cosmooth}{Predicting neurons the model
never saw, LFADS against GPFA at the same latent dimensionality. \emph{Left},
paired by session; LFADS is higher in 8 of 8. \emph{Right}, the same data against
the line of equality. Both models are scored against an identical spike target
through an identical readout, so the comparison isolates the quality of the
latent state.}

\subsection{The latent state predicts the actual lever movement}\label{sec:lever}

The model was fit to spikes. It never saw the lever. Yet a linear readout of the
generator state recovers the continuous lever position on held-out trials, at
$R^2 = \mathbf{0.47}$ to $\mathbf{0.77}$ across four sessions
(Figs.~\ref{fig:levertr} and \ref{fig:lever}).

The comparison that matters is the grey bar: predict every held-out trial using
the training-set average pull. That scores 0.26 to 0.68, because the movement is
stereotyped and the average is already a good guess. The neural readout beats it
by 0.07 to 0.21, and that margin is the part that is specific to the individual
trial rather than common to all of them.

\rfig[0.52]{lever_decode_traces.png}{levertr}{Five held-out trials from one
session. Grey is the real lever position, green is the prediction from a window
of generator state. Each trial was excluded from the fold that predicts it. The
prediction tracks onset, duration and return to rest.}

\rfig[0.95]{lever_decode.png}{lever}{Lever decoding across all four sessions.
\emph{Left}, six held-out trials, real against predicted. \emph{Right}, every
session. Grey bars predict each held-out trial with the training-set average
pull, which is a strong baseline because the movement is stereotyped. Green is
the neural readout. The readout uses a $-60$ to $+160$~ms window of generator
state, 5-fold cross-validated across trials with an inner loop for the ridge
penalty. Successful trials only.}

\subsection{Reconstruction error picks the wrong model}\label{sec:sel}

This is worth stating plainly because it nearly derailed the analysis.

The model with the best validation reconstruction, \texttt{cfg01}, turns out to
be almost useless for studying dynamics. Run its generator forward with the
inputs switched off and it recovers only \textbf{47\%} of the trajectory in mouse
1 and \textbf{46\%} in mouse 2. Most of what looked like dynamics was actually
being injected through the inferred inputs. \texttt{cfg06} recovers \textbf{97\%}
and \textbf{87\%} respectively, at a cost of 2.9\% and 4.6\% worse reconstruction.

Fig.~\ref{fig:modelsel} shows why this happens. Reconstruction error and autonomy
are close to unrelated across the fitted models, and \texttt{cfg01} sits alone in
the bottom-left corner of both panels: best reconstruction, worst autonomy. The
same pattern appears in both animals independently.

The mechanism is simple. The two configurations differ in exactly one setting:
the penalty on the controller's outputs (\texttt{kl\_co\_scale}) is $10\times$
weaker in \texttt{cfg01}. A weaker penalty lets the model explain the data by
injecting input, and reconstruction error rewards that exactly as much as
explaining it with the learned rule.

\begin{quote}
\textbf{If the claim is about intrinsic dynamics, reconstruction error is not a
valid way to choose a model.}
\end{quote}

Every result in this report uses \texttt{cfg06}. We did not simply take the most
autonomous model available (two reach 99\%), because at that level the inferred
inputs carry no signal at all and the result in \S\ref{sec:cue} disappears.

\rfig[0.98]{model_selection.png}{modelsel}{Reconstruction error and autonomy are
close to unrelated. Each point is one fitted model. Horizontal axis is validation
reconstruction loss, lower being better. Vertical axis is how much of the
trajectory the generator reproduces with all inferred input removed. The
best-reconstructing model (red square) has the worst autonomy in both animals.
The selected model (blue circle) gives up a little reconstruction for a
large gain in autonomy.}

\subsection{The rule alone regenerates single trials}

Give the generator one trial's starting state, switch off every input, and let it
run. It reproduces 94\% of that trial's latent trajectory before the cue and 82\%
after (Fig.~\ref{fig:rollout}).

Three controls say what this does and does not mean:

\begin{itemize}
\item A predictor that just uses the \textbf{condition average} scores
$R^2 = 0.02$ and $0.03$. So the trials are genuinely different from one another,
and the rollout is capturing that difference rather than a common shape.
\item \textbf{Held-out trials} score 0.96 and 0.78, against 0.93 and 0.82 on
training trials. No overfitting.
\item Feeding in \textbf{another trial's} starting state gives $R^2 < 0$.
\end{itemize}

One honest caveat. The starting state is inferred by an encoder that reads the
whole trial, so this is a statement about \emph{compressibility}, not prediction.
What it shows is that a 3-second, 15-dimensional, trial-specific trajectory
collapses to one point plus one shared rule. The rule, not the input, carries the
trajectory.

The one place it fails is right at the cue, which is exactly where something
arrives from outside. That is taken up in \S\ref{sec:cue}.

\rfig[0.95]{fig2_autonomous_rollout.png}{rollout}{The rule regenerates single
trials. Given only a trial's inferred starting state, with all external input set
to zero, the generator runs forward 150 steps. \emph{Left}, one trial: the
trajectory the full model inferred (black) against the one generated from the
starting state alone (green). \emph{Right}, all trials, as variance explained
over time, with the wrong-starting-state null shown for comparison.}

\subsection{What the learned rule looks like}\label{sec:flow}

Fig.~\ref{fig:flow} shows the fitted rule directly rather than describing it.
Each arrow is one real step of the trained generator, taken from the average
recorded state in that region, so the picture describes what the model does where
the data actually lives. A successful trial and a failed trial start close
together and end up in different places. The same structure appears in both
animals, from models fitted separately.

\rfig[0.98]{flowfield.png}{flow}{The rule the population follows, measured in both
animals. Arrows show one step of the fitted generator from the average real state
in that region, projected onto the top two latent dimensions (36\% and 59\% of
state variance). Colour is how fast the state is changing. Arrows fade where
fewer recordings pass through, so pale regions are the least certain. One hit
trial and one miss trial are overlaid.}

\subsection{The circuit holds the plan on a continuous attractor}
\label{sec:attractor}

All 512 starting points converged to \emph{different} fixed points. That is
already informative: a handful of isolated attractors would have pulled every
search to the same few places. A whole surface of them is what you get from a
continuous attractor.

Pushing on each one gives the picture in Fig.~\ref{fig:attr}. Across 60 sampled
fixed points there are, per fixed point, \textbf{33.2 directions where a push
stays}, \textbf{94.7 where it decays} with a median time constant of 84~ms, and
\textbf{0.08 where it grows}, which is to say essentially none.

No model we examined had unstable directions. The structure held in \textbf{five
separately trained models per animal} (median 34 marginal directions in mouse 1,
range 20 to 46; median 36 in mouse 2, range 8 to 50). How many marginal
directions there are depends on regularisation and mildly on sampling; that there
are some, and that nothing is unstable, does not.

This resolves the problem from the introduction without any compromise being
struck. The plan is held along the marginal directions. Noise is removed across
the stable ones. Nothing has to be traded off, because the two jobs are done by
different directions in the same space.

\rfig[0.98]{fig3_attractor_spectrum.png}{attr}{The trained network implements a
continuous attractor. \emph{Left}, exact eigenvalues at the fixed points against
the unit circle. The spectrum crowds the circle from inside with essentially
nothing outside it. \emph{Centre}, sorted eigenvalue magnitude per fixed point
(mean $\pm$ s.d.): a plateau of marginal modes near $|\lambda| = 1$ gives way to a
decaying stable spectrum. \emph{Right}, the surface of fixed points is itself
low-dimensional, comparable to the data. Model \texttt{cfg06},
\texttt{61199SOM}, 512 fixed points, 60 sampled.}

\subsection{The same structure appears in every recording and in both animals}
\label{sec:gen}

Mean successful-trial paths were compared in the model's full 15-dimensional
factor space, using ordinary Procrustes alignment, which allows only a shift, a
uniform rescale and a rotation. Nothing is bent or warped to make the curves
agree.

Across four sessions of one animal, the shape difference is 0.02 to 0.21. Across
the \emph{two animals}, fitted by completely separate models with no shared
parameters, it is \textbf{0.11}. Aligning to a time-scrambled version of the
target gives $0.99 \pm 0.01$, so 1.00 is what unrelated shapes look like. Every
comparison reaches $p < 0.001$ against that null
(Figs.~\ref{fig:crossanimal} and \ref{fig:gen}).

\rfig[0.52]{crossanimal.png}{crossanimal}{Two animals, two separately trained
models, no shared parameters. After only a shift, scale and rotation, the mean
successful-trial paths lie on top of each other. Dots mark cue onset.}

\rfig[0.98]{generalization.png}{gen}{The same trajectory and the same dynamics in
every recording. \textbf{A}, mean successful-trial path from four sessions of one
animal, each rigidly aligned to the first. \textbf{B}, the two animals, fitted by
separate models. \textbf{C}, marginally stable directions found by five
separately trained models per animal. Comparisons are made in the full
15-dimensional space; a shape difference of 1.00 is what aligning to a
time-scrambled target gives.}

\subsection{The cue shows up as a push from outside, and nobody told the model}
\label{sec:cue}

The model saw spike trains and nothing else. It was never told when the tone
played or which trials the mouse got right.

The external input it infers is flat for the whole second before the cue, then
rises sharply, peaking at \textbf{1.73$\times$ baseline at $+100$~ms} in mouse 1
and \textbf{2.18$\times$ at $+140$~ms} in mouse 2 (Fig.~\ref{fig:cue}). That is
the model discovering the cue on its own, in two animals separately.

The push is also larger on trials the mouse got right
($P(\text{Hit} > \text{Miss}) = 0.619$, $p = 1.8 \times 10^{-10}$; and $0.813$,
$p = 9.0 \times 10^{-60}$, against a 2000-fold label-shuffle null whose 95th
percentile is 0.532). How much larger differs between the two animals; that there
is a cue-locked push at all does not.

\rfig[0.95]{fig4_cue_input.png}{cue}{The cue arrives as an inferred input. Size of
the inferred external input over time, split by whether the mouse succeeded, for
both animals. The model was given spikes only and was never told when the cue
occurred or which trials succeeded. Shading is the 95\% confidence interval on
the mean.}

\subsection{Where you push determines whether it sticks}\label{sec:steer}

Push the full nonlinear network along a marginal direction, let go, and
\textbf{75\%} of the displacement is still there 1.2~s later. Push it in a
typical other direction and \textbf{8\%} survives (Fig.~\ref{fig:steer}). A
one-off intervention therefore only persists if it is aimed within the manifold.

That raises the question that actually matters for intervention: is the direction
we care about, the one separating success from failure, one of the cheap ones?
Partly. It has 0.352 of its length inside the manifold against 0.254 expected by
chance, a factor of $\mathbf{1.38}$, and it beat all 400 random directions tested
($p < 0.0025$).

This should not be oversold. Most of the behavioural direction is still off the
manifold, which means steering behaviour would need sustained input rather than a
single pulse. But the geometry is not arbitrary with respect to behaviour, and
that is what makes aiming worth attempting.

\rfig[0.95]{fig5_steering_cost.png}{steer}{The cost of intervening depends on
direction. \textbf{A}, the full nonlinear generator is displaced from a fixed
point and released. A push along a marginal direction keeps 75\% of its size
after 1.2~s; a typical other direction decays to 8\%. Shading is s.d.\ across
states. \textbf{B}, how much of the success-versus-failure direction lies inside
the manifold (orange) against 400 random directions (grey).}

\subsection{Why none of the above can run live}

Every estimate so far uses an encoder that reads the whole trial, including what
happens after the moment being estimated. That is fine for understanding and
useless for intervening. A live detector has to commit on the rising edge of the
signal, with no access to the peak or to anything later.

Fig.~\ref{fig:pipeline} shows the causal front end built for this. Every stage
looks only backwards in time, and the filter and baseline state carry across file
boundaries and are never reset, which is why replaying a recording produces
output identical to the live stream.

\rfig[0.95]{fig6_causal_pipeline.png}{pipeline}{Turning raw signal into alarms.
Median subtraction across channels, a causal 300~Hz high-pass, a per-channel
adaptive threshold at $4\times$ a running MAD estimate, threshold crossings
accumulated into 20~ms bins, z-scoring, and the last 12 bins stacked into a
768-number vector projected onto one direction. An alarm fires when the score
stays above threshold for two bins.}

\subsection{The detector flags movement before it happens}

On held-out data the detector reached \textbf{AUC 0.877}, catching 60\% of
movements at about 5 false alarms per minute of rest, and it never sees the cue.
Alarms preceded lever motion by a median of \textbf{63~ms}, with \textbf{92\%}
arriving before the movement ($n = 199$). The unsupervised CUSUM baseline on the
same features did worse, which tells us the learned projection is doing the work,
not the front end.

Running live, streaming from the amplifier and calibrating itself from the lever
as the session went on, it flagged movement a median of \textbf{80~ms} before the
lever moved, with \textbf{90\%} of alarms early and \textbf{71\%} followed by a
real movement (Fig.~\ref{fig:live}). Of the alarms that caught a movement,
\textbf{46 preceded self-initiated pulls} and \textbf{20 preceded cued pulls},
from one template that was never told which kind it was looking at.

\paragraph{Speed.} Median compute was \textbf{under 1~ms per 20~ms bin} on a
single CPU core, under 5\% of the available budget. So the remaining delay is
structural, not computational: the 20~ms bin plus the 40~ms confirmation window
give a 60~ms floor. This lines up with the independent analysis in
\S\ref{sec:crossday}, which finds movement becomes decodable about 140~ms before
onset. 140 minus 60 predicts about 80~ms of usable lead, and 80~ms is what the
live system delivers. The detector is running at the limit set by its own bin
size rather than by its statistics.

\rfig[0.95]{fig7_live_results.png}{live}{Live performance, streaming from the
amplifier and calibrating from the lever as the session ran. \textbf{A}, how far
ahead each alarm was; green for alarms before the lever moved (90\%), red for
after, dashed line at the median of $+80$~ms. \textbf{B}, what happened after
every alarm: 20 preceded cued pulls, 46 preceded self-initiated pulls, 3 fired
late, and 28 were followed by nothing (71\% precision).}

\subsection{Bounded scores behave better live, and offline metrics cannot see it}
\label{sec:readout}

The Fisher discriminant and the PCA-whitened cosine scored the same offline
(AUC 0.83 to 0.85). They behave differently in the tails. The discriminant is
unbounded, and during rest it produced excursions of 6 to 23 standard deviations.
Live, that meant it kept drifting across its threshold and firing on rare
artefacts. The cosine is capped at $[-1,1]$ by construction, with a rest tail
near 0.6, so a strange input cannot produce an arbitrarily large score.

\begin{quote}
\textbf{At equal accuracy, prefer the bounded score for anything that runs live.}
Offline AUC cannot see this difference at all.
\end{quote}

\subsection{Correcting for how often the animal moves}

The mouse pulls the lever about every 3.7~s. A detector firing at random would
therefore land near a movement fairly often, which makes raw hit rate
uninterpretable.

Against a circular-shift null with 5000 permutations, the detector achieved
precision \textbf{74.2\% against 39.3\% $\pm$ 5.2 expected} ($z = +6.7$) and
recall \textbf{35.3\% against 16.8\% $\pm$ 2.3} ($z = +8.2$), beating 99.9\% and
100\% of permutations (Fig.~\ref{fig:shuffle}). Lead times were
\textbf{6.3$\times$ more concentrated} in the 0 to 150~ms window than a uniform
distribution would give (62.7\% against 10\%).

Note that chance precision here is 39\%, not zero. Without this control the
result would have looked far better than it is.

\rfig[0.85]{fig8_shuffle_control.png}{shuffle}{Correcting for the base rate.
Alarm times are rotated rigidly around the session clock across 5000
permutations, which keeps the alarm rate and any rhythm and destroys only the
alignment to movement. Grey is the null distribution, orange the measured value.}

\subsection{Across days the coordinates rotate, the dynamics do not}
\label{sec:crossday}

Applying a model to the next day's recording without any adjustment collapses it
to \textbf{AUC 0.485}, which is chance, with false alarms rising to 21.8 per
minute.

After a rigid Procrustes alignment in a 15-dimensional latent space, using
pull-triggered averages as landmarks, the same model reaches \textbf{AUC 0.740},
against a same-day ceiling of 0.768. So the information is still there. What
changed is the coordinate system the electrodes report it in, which fits reports
of stable latent dynamics underneath unstable single-channel readouts
\cite{gallego2020,degenhart2020}.

Seeding the calibration with the previous day's data nearly doubled early-session
recall (32\% to 56\% at 10\% of a session), with the advantage closing by 30\% as
live data accumulated. Separately, time-resolved decoding of movement against
matched rest becomes reliable \textbf{140 to 260~ms before onset} at a peak AUC
of 0.85. That number sets the ceiling on how early any causal system could fire.

\subsection{What the second recording area contributes}

The rig already carries two probes, so the question we asked of the premotor
recording was not whether it buys extra lead time, but whether M1 and M2 together
can \emph{tell self-initiated movements from cued ones}.

On lead time it adds nothing. The two population signals cross-correlate at
\textbf{0~ms lag} ($r = 0.98$ to $0.99$ on both days), their ramps have
indistinguishable latencies, and using both probes detected no earlier than M1
alone ($-120$ against $-140$~ms). This is consistent with
ongoing, unpublished work in this laboratory on how motor areas are
coordinated.

M2 did look better at first, 75.7\% against 71.9\% balanced accuracy, but that
advantage failed every control we applied. It reversed for long cue-to-movement
delays ($+2.6$ to $-4.0$ points), vanished when labels were validated against
reward (81.8 against 81.5), and was absent altogether for genuinely
self-initiated movements ($n = 447$, not significant).

That pattern places M2's extra signal in the \textbf{cue-evoked response}. Which
is not an earlier motor intention signal, but is precisely the component that
distinguishes cued from self-initiated trials, and therefore is the answer to the
question we actually asked.

\section{Discussion}

The circuit holds a plan by geometry rather than by tuning. Information sits along
a surface of marginally stable states, and noise is removed in the directions
leading off it. Nothing has to be balanced against anything else, because the two
requirements are met by different directions in the same space. The structure
survives retraining, survives changing the regularisation, and appears in two
animals whose models share no parameters. The same model also finds the cue on
its own, having never been told it exists.

This connects motor preparation to a mechanism already established for head
direction and for holding the eyes still
\cite{seung1996,zhang1996,chaudhuri2019}, which suggests continuous attractors
may be a general solution wherever a cortical circuit has to hold a value that
can change.

\paragraph{What this means for intervening.} It is a constraint, not a recipe. A
single push only lasts if it is aimed within the manifold, 75\% against 8\%. The
behaviourally useful direction leans toward the manifold but is nowhere near
contained by it, so any real controller will need sustained input with a
substantial component in the expensive directions.

The failed LQR analysis is worth keeping in view rather than hiding. The model's
own four inferred input channels turned out to have essentially no authority over
a 128-dimensional state, which is itself informative: a physical intervention has
to reach the state more directly than the model's input parameterisation
suggests. That is an argument for the nanoscale and optical interfaces developed
in this laboratory \cite{jayant2019,gonzales2025}, and it distinguishes this
problem from conventional decoding prostheses \cite{gilja2012}, where the aim is
to read the state rather than to move it.

\paragraph{What Part B settles and what it does not.} It settles the timing.
Movement intention is readable causally, from unsorted multiunit activity, about
80~ms before the movement, with compute to spare, and the result survives a null
that corrects for how often the animal moves. It also identifies the real
engineering obstacle, which is that the readout rotates between days, and shows
that fixing it is cheap.

It does \textbf{not} settle state estimation. The detector is discriminative. It
answers ``is a movement coming?'' and not ``where on the manifold is the circuit
right now?''. Aiming a push, which is what \S\ref{sec:steer} calls for, needs the
second question answered. Part~B should be read as establishing that the timing
and the front end are feasible, leaving causal estimation of the manifold
coordinate as the substantive open problem.

\paragraph{Limitations.} Four matter. The Part~A estimates are \textbf{acausal},
because the encoder reads the whole trial; nothing in Part~A could be executed
online as it stands. The Part~B detector is discriminative rather than
generative. The number of marginal directions depends on regularisation and on
sampling, ranging from 8 to 50 across models, so the count should not be
over-interpreted even though its existence is robust. And the sample is small:
the live validation rests on one session in one animal, and the cross-area
comparison on two sessions in one animal. The controls are strong; the number of
animals is not.

\section{Conclusions and Future Directions}

Fitting a dynamical model to mouse motor cortex and then examining the rule it
learned shows that the circuit holds a movement plan on a continuous attractor:
roughly 33 directions in which a push stays put, roughly 95 in which it is
corrected within 84~ms, and essentially none in which it runs away. The same
latent state predicts neurons the model never saw, predicts the lever movement it
was never shown, and regenerates individual trials from a single point plus one
shared rule. The geometry is not incidental to behaviour: the direction
separating success from failure sits preferentially, though only modestly, in the
cheap subspace.

Separately, movement intention can be read causally and in real time from
unsorted multiunit activity, 80~ms before the lever moves, at a compute cost far
below the real-time budget, and the result survives a null that corrects for how
often the animal moves anyway.

The gap between these two halves is precise and is the obvious next step. Part~A
gives a target and a price. Part~B gives a trigger. Closing the loop needs a
\emph{coordinate}, not an alarm.

\paragraph{Next steps.}
\begin{enumerate}
\item \textbf{Causal state estimation.} Replace the bidirectional encoder with a
forward-only one over a fixed window of the past, so the system outputs a
position on the manifold rather than a binary alarm. This is the direct bridge
between the two halves of this work.
\item \textbf{The experiment that could prove this wrong.} Deliver optogenetic
stimulation along the manifold and across it. \S\ref{sec:steer} predicts 75\%
retention against 8\%. If the two conditions behave the same, the picture in this
report is wrong. The validated detector from Part~B is already sufficient to
trigger such stimulation at the right moment.
\item \textbf{Learning.} Whether the number of marginal directions changes as an
animal learns the task, which is testable in the multi-day recordings that
already exist.
\item \textbf{Method.} Treating autonomy as a routine selection criterion
whenever LFADS-class models are used to make claims about dynamics rather than to
denoise.
\end{enumerate}

\section*{Acknowledgements}

I thank \textbf{Om Kolhe} for day-to-day mentorship and guidance throughout this
project, and for the dual-probe recordings analysed here; \textbf{Hammad Khan}
and \textbf{Shivam} for many helpful discussions and technical guidance; and
\textbf{Prof.\ Krishna Jayant} for hosting this work in the Nano Neurotechnology
Laboratory and for direction on its scope. All electrophysiological recordings
analysed in this report were performed by members of the laboratory; my
contribution was computational. Thanks also to graduate advisor \textbf{Monish
Lokhande}. Supported by the Purdue Summer Undergraduate Research Fellowship
(SURF).

\begin{thebibliography}{99}

\bibitem{tanji1976} Tanji, J. \& Evarts, E. V. (1976). Anticipatory activity of
motor cortex neurons in relation to direction of an intended movement.
\emph{J. Neurophysiol.} 39(5), 1062--1068.

\bibitem{churchland2012} Churchland, M. M., Cunningham, J. P., Kaufman, M. T., et
al. (2012). Neural population dynamics during reaching. \emph{Nature} 487,
51--56.

\bibitem{cunningham2014} Cunningham, J. P. \& Yu, B. M. (2014). Dimensionality
reduction for large-scale neural recordings. \emph{Nat. Neurosci.} 17,
1500--1509.

\bibitem{vyas2020} Vyas, S., Golub, M. D., Sussillo, D. \& Shenoy, K. V. (2020).
Computation through neural population dynamics. \emph{Annu. Rev. Neurosci.} 43,
249--275.

\bibitem{seung1996} Seung, H. S. (1996). How the brain keeps the eyes still.
\emph{PNAS} 93(23), 13339--13344.

\bibitem{zhang1996} Zhang, K. (1996). Representation of spatial orientation by
the intrinsic dynamics of the head-direction cell ensemble.
\emph{J. Neurosci.} 16(6), 2112--2126.

\bibitem{chaudhuri2019} Chaudhuri, R., Ger\c{c}ek, B., Pandey, B., Peyrache, A.
\& Fiete, I. (2019). The intrinsic attractor manifold and population dynamics of
a canonical cognitive circuit across waking and sleep. \emph{Nat. Neurosci.} 22,
1512--1520.

\bibitem{pandarinath2018} Pandarinath, C., O'Shea, D. J., Collins, J., et al.
(2018). Inferring single-trial neural population dynamics using sequential
auto-encoders. \emph{Nat. Methods} 15, 805--815.

\bibitem{sussillo2013} Sussillo, D. \& Barak, O. (2013). Opening the black box:
low-dimensional dynamics in high-dimensional recurrent neural networks.
\emph{Neural Comput.} 25(3), 626--649.

\bibitem{golub2018} Golub, M. D. \& Sussillo, D. (2018). FixedPointFinder: a
Tensorflow toolbox for identifying and characterizing fixed points in recurrent
neural networks. \emph{J. Open Source Softw.} 3(31), 1003.

\bibitem{sedler2023} Sedler, A. R. \& Pandarinath, C. (2023). lfads-torch: a
modular and extensible implementation of latent factor analysis via dynamical
systems. \emph{arXiv:2309.01230}.

\bibitem{yu2009} Yu, B. M., Cunningham, J. P., Santhanam, G., Ryu, S. I., Shenoy,
K. V. \& Sahani, M. (2009). Gaussian-process factor analysis for low-dimensional
single-trial analysis of neural population activity. \emph{J. Neurophysiol.}
102(1), 614--635.

\bibitem{pei2021} Pei, F., Ye, J., Zoltowski, D., et al. (2021). Neural Latents
Benchmark '21: evaluating latent variable models of neural population activity.
\emph{arXiv:2109.04463}.

\bibitem{keshtkaran2022} Keshtkaran, M. R., Sedler, A. R., Chowdhury, R. H., et
al. (2022). A large-scale neural network training framework for generalized
estimation of single-trial population dynamics. \emph{Nat. Methods} 19,
1572--1577.

\bibitem{trautmann2019} Trautmann, E. M., Stavisky, S. D., Lahiri, S., et al.
(2019). Accurate estimation of neural population dynamics without spike sorting.
\emph{Neuron} 103(2), 292--308.

\bibitem{gallego2020} Gallego, J. A., Perich, M. G., Chowdhury, R. H., Solla, S.
A. \& Miller, L. E. (2020). Long-term stability of cortical population dynamics
underlying consistent behavior. \emph{Nat. Neurosci.} 23, 260--270.

\bibitem{degenhart2020} Degenhart, A. D., Bishop, W. E., Oby, E. R., et al.
(2020). Stabilization of a brain--computer interface via the alignment of
low-dimensional spaces of neural activity. \emph{Nat. Biomed. Eng.} 4, 672--685.

\bibitem{page1954} Page, E. S. (1954). Continuous inspection schemes.
\emph{Biometrika} 41(1/2), 100--115.

\bibitem{fisher1936} Fisher, R. A. (1936). The use of multiple measurements in
taxonomic problems. \emph{Ann. Eugen.} 7(2), 179--188.

\bibitem{gilja2012} Gilja, V., Nuyujukian, P., Chestek, C. A., et al. (2012). A
high-performance neural prosthesis enabled by control algorithm design.
\emph{Nat. Neurosci.} 15, 1752--1757.

\bibitem{gonzales2025} Gonzales, D. L., Khan, H. F., Keri, H. V. S., Yadav, S.,
Steward, C., Muller, L. E., Pluta, S. R. \& Jayant, K. (2025). A translaminar
spacetime code supports touch-evoked traveling waves. \emph{Science Advances} 11,
eadr4038.

\bibitem{jayant2019} Jayant, K., Wenzel, M., Bando, Y., et al. (2019). Flexible
nanopipettes for minimally-invasive intracellular electrophysiology in vivo.
\emph{Cell Reports} 26(1), 266--278.

\end{thebibliography}

\end{document}
